%% ==========================================================================
%%  Approach 9 --- the conditioned-remainder induction.  OPEN.
%%
%%  Source: CRI.md (2026-08-09) and updates/update_47/.  Every lemma here is proved and
%%  machine-checked.  conj:cri-depth, once the open statement here, is REFUTED
%%  (prop:cri-depth-false, 2026-08-10).  The surviving route is the
%%  witness-carrying invariant of sec:cri-witness, whose one open clause is
%%  conj:cri-progress.
%%
%%  This is Approach 3 done as an INDUCTION rather than as a search.  The peel
%%  frame of Approach 3 is retained wholesale --- nothing there is retracted ---
%%  but the state space is replaced by one whose elements are partial
%%  allocations under contracted cost functions.
%%
%%  Scripts: updates/update_47/cri_anchor.py    (the lemmas)
%%           updates/update_47/cri_sweep.py     (0 bad roots; the refutations)
%%           updates/update_47/cri_stuck.py     (the contingency tables)
%%           updates/update_47/cri_where.py     (stuckness sits at |R| <= 2)
%%           updates/update_47/cri_lookahead.py (no dead state at |R| >= 3)
%%           updates/update_47/cri_witnesses.py (the minimal witness for each refutation)
%% ==========================================================================
\section{Approach 9: The Conditioned-Remainder Induction}
\label{sec:approach9}

\Cref{sec:approach3} left exactly one statement open,
Conjecture~\ref{conj:balance-rule}, and \Cref{rem:testing-exhausted} recorded
that random testing had run out of power on it. What followed narrowed the
target without closing it: three families of candidate invariant were
eliminated --- a local sufficient condition on the peel
(Proposition~\ref{prop:lemmas-insufficient}), a predicate on the pair
(state, move) (\Cref{rem:no-separating-feature}), and the commitment predicate
on the state alone (\Cref{rem:commit-predicate}) --- leaving no candidate at
all.

This section changes the state space rather than hunting for another predicate
on the old one. The diagnosis is that a workload profile is the wrong object: it
is an arbitrary tuple of subsets, there are $(2^m)^n$ of them, and an individual
profile means nothing. What follows replaces it with a space whose elements are
\emph{partial allocations}, and in which the peel frame's known dead ends
provably do not lie.

% --------------------------------------------------------------------------
\subsection{The frame}
\label{sec:cri-frame}

\begin{definition}[CR state]
\label{def:crstate}
A \emph{conditioned-remainder state}, or CR state, is a pair $(\alloc, R)$ with
$R \subseteq \items$ and $\alloc = (\bundle{1},\dots,\bundle{n})$ a partition of
$D := \items \setminus R$. Its \emph{profile} is
\[
  \workprofile_i \;:=\; \bundle{i} \cup R,
\]
so every agent is still on the hook for the whole undecided remainder, and its
\emph{contracted cost functions} are
\[
  \cost^{R}_i(T) \;:=\; \cost_i(T \cup R) - \cost_i(R),
  \qquad T \subseteq D .
\]
The moves are \emph{assignment} --- give some $a \in R$ to some agent $x$,
setting $\bundle{x}' = \bundle{x} \cup \set{a}$ and $R' = R \setminus \set{a}$
--- and \emph{relabelling}, $\alloc \mapsto \alloc \circ \sigma$. The
\emph{root} has $\alloc$ empty and $R = \items$; a state is \emph{terminal} when
$R = \emptyset$; and \emph{legality} is the invariant of \Cref{def:peel}: no
positive-weight cycle and $\pathw{\workprofile}(i) \le 1$ for every $i$.
\end{definition}

In the language of \Cref{def:peel} an assignment is one \emph{atomic block}:
relieve all $n-1$ non-owners of $a$ at once. So a CR state is a peel state, a CR
move is $n-1$ consecutive peels, and the frame differs from
Conjecture~\ref{conj:h1pp} exactly by dropping the legality requirement at the
$n-2$ intermediates \emph{within} a block.

\begin{conjecture}[CRI --- \textbf{open}]
\label{conj:cri}
From the root, some sequence of assignments and relabellings reaches a terminal
state with every intermediate CR state legal.
\end{conjecture}

Conjecture~\ref{conj:h1pp} implies Conjecture~\ref{conj:cri} implies
Conjecture~\ref{conj:main}, the last because a terminal CR state \emph{is} an
allocation and its legality is exactly $\max_i \pathw{\alloc}(i) \le 1$. So
Conjecture~\ref{conj:cri} is strictly the weaker target of the two peel-frame
statements, and unlike Conjecture~\ref{conj:balance-rule} it is not obviously
stronger than what is wanted.

% --------------------------------------------------------------------------
\subsection{Why the frame is different}
\label{sec:cri-different}

Three things separate it from a restriction of the peel process.

\paragraph{The space is small.} CR states are partial functions $\items \to
\agents$, so there are $(n+1)^m$ of them against $(2^m)^n$ profiles: at
$n = 3, m = 6$, $4{,}096$ against $262{,}144$. Reachability is therefore
decidable exhaustively well past every sweep of \Cref{sec:approach3}.

\paragraph{The known dead ends are not CR states.} A profile is a CR state if
and only if the sets $\workprofile_i \setminus \bigcap_k \workprofile_k$ are
pairwise disjoint and cover $\items \setminus \bigcap_k \workprofile_k$. Both
witnesses recorded in \Cref{sec:deadends} repeat an item across bundles ---
$(\set{a_1,a_2},\set{g},\set{g})$ and
$(\set{g_2},\set{g_2},\set{g_1,g_3,g_4})$ --- so neither is a CR state. The
obstruction of Proposition~\ref{prop:deadends} is removed rather than worked
around.

\paragraph{The induction is restored.} This is the substantive point, and it is
Lemmas~\ref{lem:cri-contract} and~\ref{lem:cri-witness} below: a CR state is a
\emph{solved instance of the problem on a contracted ground set}, and a move
un-contracts one element. Conjecture~\ref{conj:cri} is therefore an induction in
which the instance \emph{grows} while a witness is maintained --- the shape of
the proof of~\cite{BKNS22} --- run on the conditioned cost functions that
dissolve the obstruction of Theorem~\ref{thm:obstruction}.
\Cref{sec:approach1} grew the instance with the \emph{unconditioned} costs and
died on that obstruction; \Cref{sec:approach3} identified the right conditioning
in \Cref{rem:conditioned} and then abandoned the induction for a state-space
search. This section is the synthesis of the two.

% --------------------------------------------------------------------------
\subsection{The calculus}
\label{sec:cri-calculus}

\begin{lemma}[Contraction stays in the class]
\label{lem:cri-contract}
$\cost^{R}_i$ is normalised with every marginal in $\set{0,1}$. Hence
$(D, \cost^{R})$ is again a negative dichotomous instance.
\end{lemma}

\begin{proof}
$\cost^{R}_i(\emptyset) = \cost_i(R) - \cost_i(R) = 0$, and for
$g \notin T \subseteq D$,
\[
  \cost^{R}_i(T \cup \set{g}) - \cost^{R}_i(T)
  \;=\; \cost_i\big((T \cup R) \cup \set{g}\big) - \cost_i(T \cup R)
  \;\in\; \set{0,1}
\]
because $\cost_i$ is dichotomous (\Cref{def:dichcost}).
\end{proof}

\begin{lemma}[A CR state is a witness on the contracted instance]
\label{lem:cri-witness}
$\edgew{\workprofile}(i,k) = \cost^{R}_i(\bundle{i}) - \cost^{R}_i(\bundle{k})$.
Consequently $(\alloc,R)$ is legal if and only if $\alloc$ witnesses
Conjecture~\ref{conj:main} on the instance $(D, \cost^{R})$.
\end{lemma}

\begin{proof}
$\edgew{\workprofile}(i,k) = \cost_i(\bundle{i} \cup R) - \cost_i(\bundle{k}
\cup R) = \big[\cost^{R}_i(\bundle{i}) + \cost_i(R)\big] -
\big[\cost^{R}_i(\bundle{k}) + \cost_i(R)\big]$, and the $\cost_i(R)$ cancel.
The two envy graphs are therefore equal arc for arc, so the legality conditions
coincide.
\end{proof}

\begin{lemma}[CR arc update]
\label{lem:cri-arcupdate}
Let $(\alloc,R) \to (\alloc',R')$ assign $a$ to $x$, and write
$\beta_i(T) := \cost^{R'}_i(T \cup \set{a}) - \cost^{R'}_i(T) \in \set{0,1}$ for
the marginal of $a$ to agent $i$ on $T$, conditioned on $R'$. Then for $i \ne k$,
\[
  \edgew{\alloc'}(i,k) - \edgew{\alloc}(i,k)
  \;=\; \beta_i(\bundle{k})\,[k \ne x] \;-\; \beta_i(\bundle{i})\,[i \ne x].
\]
So arcs \emph{into} $x$ fall by $\beta_i(\bundle{i})$, arcs \emph{out of} $x$
rise by $\beta_x(\bundle{k})$, and arcs between third parties move by
$\beta_i(\bundle{k}) - \beta_i(\bundle{i}) \in \set{-1,0,1}$.
\end{lemma}

\begin{proof}
First the contraction identity: for $T \subseteq D$,
\begin{align*}
  \cost^{R'}_i(T \cup \set{a}) - \cost^{R'}_i(\set{a})
  &= \big[\cost_i(T \cup \set{a} \cup R') - \cost_i(R')\big]
   - \big[\cost_i(\set{a} \cup R') - \cost_i(R')\big] \\
  &= \cost_i(T \cup R) - \cost_i(R) \;=\; \cost^{R}_i(T).
\end{align*}
Applying it to both terms of $\edgew{\alloc}(i,k)$, the
$\cost^{R'}_i(\set{a})$ cancel and
\[
  \edgew{\alloc}(i,k) = \cost^{R'}_i(\bundle{i} \cup \set{a})
  - \cost^{R'}_i(\bundle{k} \cup \set{a})
  = \big[\cost^{R'}_i(\bundle{i}) + \beta_i(\bundle{i})\big]
  - \big[\cost^{R'}_i(\bundle{k}) + \beta_i(\bundle{k})\big].
\]
Meanwhile $\edgew{\alloc'}(i,k) = \cost^{R'}_i(\bundle{i}') -
\cost^{R'}_i(\bundle{k}')$ with $\bundle{x}' = \bundle{x} \cup \set{a}$ and
$\bundle{k}' = \bundle{k}$ for $k \ne x$. Subtracting in each of the three cases
$i,k \ne x$; $k = x$; $i = x$ gives the stated formula.
\end{proof}

\begin{corollary}[The additive collapse]
\label{cor:cri-additive}
If the $\cost_i$ are additive then $\beta_i$ is constant, so third-party arcs do
not move and only the arcs at $x$ change.
\end{corollary}

The corollary localises the difficulty sharply: \emph{the entire obstruction to
the induction is the variation of $\beta_i$ across bundles}, and the additive
case is closed independently by Theorem~\ref{thm:binadd}.

\begin{proposition}[The first chore, exactly]
\label{prop:cri-first}
From the root, assigning $a$ to $x$ is legal if and only if
$\beta_x \le \beta_k$ for every $k$, where
$\beta_i := \cost_i(\items) - \cost_i(\items \setminus \set{a})$: the chore must
go to an agent of \emph{minimal} marginal.
\end{proposition}

\begin{proof}
At the root every $\bundle{i}$ is empty, so afterwards $\workprofile_x = \items$
and $\workprofile_i = \items \setminus \set{a}$ for $i \ne x$. Hence
$\edgew{}(i,k) = 0$ for $i,k \ne x$, $\edgew{}(i,x) = -\beta_i$ and
$\edgew{}(x,k) = \beta_x$. For $n \ge 3$ the first group forces $\subsidy_i = c$
for all $i \ne x$, and the remaining constraints read
$c + \beta_x \le \subsidy_x \le c + \min_{i \ne x}\beta_i$. This interval meets
$\set{0,1}$ exactly when $\beta_x \le \min_{i \ne x}\beta_i$: if $\beta_x = 0$
take $c = \subsidy_x = 0$, and if $\beta_x = 1$ then the minimum is also $1$ and
$c = 0$, $\subsidy_x = 1$ works, while $\beta_x = 1 > 0 = \min_i \beta_i$ leaves
it empty. For $n = 2$ the first group is vacuous and the same two constraints
give the same criterion. Legality is then Lemma~\ref{lem:potential-form}.
\end{proof}

This is the exact mirror of Proposition~\ref{prop:first-peel}, which takes a
chore \emph{from} an agent of maximal marginal. Such an $x$ always exists, so
the base case never blocks.

\begin{lemma}[Free assignment]
\label{lem:cri-free}
If $\beta_x(\bundle{k}) = 0$ for every $k \ne x$ and $\beta_i(\bundle{k}) \le
\beta_i(\bundle{i})$ for all $i,k \ne x$, then every arc weight weakly decreases,
so the assignment preserves legality.
\end{lemma}

\begin{proof}
Immediate from Lemma~\ref{lem:cri-arcupdate}: arcs into $x$ fall, arcs out of
$x$ rise by $\beta_x(\bundle{k}) = 0$, and third-party arcs move by
$\beta_i(\bundle{k}) - \beta_i(\bundle{i}) \le 0$. Path and cycle weights are
sums of arc weights.
\end{proof}

% --------------------------------------------------------------------------
\subsection{What the frame does not fix}
\label{sec:cri-refuted}

Conjecture~\ref{conj:cri} survives every instance tested (\Cref{sec:cri-evidence}),
but the local statements one would want to prove it with do not. One instance
carries most of the damage.

\begin{example}[A zero-subsidy dead end]
\label{ex:cri-deadend}
Take $\items = \set{a,b,c}$ with the dichotomous costs
\begin{center}\small
\begin{tabular}{@{}lcccccccc@{}}
\toprule
$S$ & $\emptyset$ & $a$ & $b$ & $c$ & $ab$ & $ac$ & $bc$ & $abc$ \\
\midrule
$\cost_1$ & 0 & 0 & 0 & 0 & 0 & 1 & 1 & 1 \\
$\cost_2$ & 0 & 0 & 1 & 1 & 1 & 1 & 1 & 1 \\
$\cost_3$ & 0 & 1 & 1 & 1 & 1 & 1 & 2 & 2 \\
\bottomrule
\end{tabular}
\end{center}
and the CR state $\alloc = (\emptyset, \set{b}, \set{a})$, $R = \set{c}$, whose
profile is $\workprofile = (\set{c},\set{bc},\set{ac})$ with
\[
  \big(\edgew{\workprofile}(i,k)\big) =
  \begin{pmatrix} 0 & -1 & -1 \\ 0 & 0 & 0 \\ 0 & -1 & 0 \end{pmatrix},
  \qquad \pathw{\workprofile} = (0,0,0).
\]
The state is reachable from the root and \emph{exactly envy-free}: it needs no
subsidy at all. Yet each of the three ways to place the last chore creates a
positive-weight cycle --- $c \to 1$ gives $\edgew{}(2,3) + \edgew{}(3,2) = 1$,
$c \to 2$ gives $\edgew{}(3,1) + \edgew{}(1,3) = 1$, and $c \to 3$ gives
$\edgew{}(2,1) + \edgew{}(1,2) = 1$ --- so no continuation is even
envy-freeable. Conjecture~\ref{conj:main} nevertheless holds on this instance:
$(\set{c},\set{ab},\emptyset)$ has an identically zero envy matrix.
\end{example}

\begin{remark}[What Example~\ref{ex:cri-deadend} rules out]
\label{rem:cri-nopotential}
It refutes three statements at once. \emph{Pointwise} extendability --- every
legal non-terminal CR state has a legal assignment --- fails, as it does in the
peel frame (Conjecture~\ref{conj:h1}); $3{,}024$ reachable states are of this
kind. The \emph{last-rung} lemma --- a legal state with one chore left can
always place it --- fails, and $1{,}765$ of the $1{,}778$ stuck states have
$\abs{R} = 1$, so this is where essentially all stuckness lives.

The third is the one that constrains a proof. The doomed state has
$\pathw{\workprofile} = (0,0,0)$: it is as good as a CR state can possibly be,
and the instance itself is not a counterexample to
Conjecture~\ref{conj:main}. \emph{So no potential function measuring the quality
of the current state can serve as the invariant}, since the state that must be
avoided is optimal. \Cref{rem:no-separating-feature} reached a similar
conclusion in the peel frame by exhausting the features of a move; here it
follows from one instance, and for a structural reason.
\end{remark}

\begin{remark}[Balance does not transplant]
\label{rem:cri-balance-fails}
The same state refutes the one idea that survived \Cref{sec:approach3}. Its
bundle sizes are $(0,1,1)$ with one chore undecided, so the completion with
sizes $(1,1,1)$ exists and the state \emph{admits a balanced terminal} in the
sense of Observation~\ref{obs:balance-characterises} --- and it is dead. Over
the corpus, $132$ dead CR states admit a balanced completion. The certificate of
Conjecture~\ref{conj:balance-rule} is therefore specific to the peel frame and
does not carry to this one.
\end{remark}

Three further statements fail. Each has an explicit minimal witness, recorded
and re-verified in \texttt{cri\_witnesses.py}.
\emph{Relabelling is load-bearing}: with
$\cost_1 = \cost_2$ charging $1$ only for the grand bundle and
$\cost_3(S) = \min(\abs{S},2)$, assignments alone cannot reach a legal terminal
from the root while one relabelling suffices, and $3$ such instances occur in
the exhaustive family. \emph{Death is not always immediate}: at $n = 3$,
$m = 5$ there is a dead state with two legal assignments available, so one-step
legality cannot certify a rule and Lemma~\ref{lem:cri-free} is sound for a move
and worthless for a schedule. And the commitment predicate of
\Cref{rem:commit-predicate} transplants \emph{backwards}: of $1{,}210$ states in
which some agent already carries conditioned cost $2$, none is dead.

% --------------------------------------------------------------------------
\subsection{Where the difficulty is, and the one open statement}
\label{sec:cri-evidence}

Every sweep below enumerates its state space completely, with no cap, and draws
from all seven generators of \texttt{updates/update\_44/counterexample\_hunt.py} rather
than the single uniform sampler used throughout \Cref{sec:approach3}.

\begin{observation}[The sweep]
\label{obs:cri-sweep}
Over the complete exhaustive $n = m = 3$ family --- all $9{,}880$ instances ---
together with $438$ adversarial instances at $n = 3$ ($m \le 7$), $n = 4$
($m \le 6$) and $n = 5$ ($m \le 5$), there is \emph{no bad root}:
Conjecture~\ref{conj:cri} holds on every one. Without the relabelling move there
are $3$ bad roots, all in the exhaustive family
(\texttt{updates/update\_47/cri\_sweep.py}).
\end{observation}

The dead states are confined to the end of the induction, and sharply so.

\begin{observation}[Death lives in the last two rungs]
\label{obs:cri-depth}
Over $312{,}436$ reachable non-terminal CR states:
\begin{center}\small
\begin{tabular}{@{}lrrr@{}}
\toprule
$\abs{R}$ & dead & reachable & dead share \\
\midrule
$1$        & $1{,}230$ & $183{,}984$ & $0.669\%$ \\
$2$        & $45$      & $93{,}909$  & $0.048\%$ \\
$\ge 3$    & $\mathbf{0}$ & $34{,}543$ & $\mathbf{0\%}$ \\
\bottomrule
\end{tabular}
\end{center}
Correspondingly, playing an arbitrary legal assignment while $\abs{R} > 3$ and
then searching the last three chores reached a terminal on all $10{,}045$
instances tested, with the play above the threshold chosen at random so that the
result is not an artefact of a tie-break
(\texttt{updates/update\_47/cri\_lookahead.py}).
\end{observation}

\begin{conjecture}[Depth --- \textbf{refuted}]
\label{conj:cri-depth}
Every reachable legal CR state with $\abs{R} \ge 3$ is live: some sequence of
assignments and relabellings from it reaches a terminal state through legal
states.
\end{conjecture}

\begin{proposition}
\label{prop:cri-depth-implies}
Conjecture~\ref{conj:cri-depth} implies Conjecture~\ref{conj:cri}, hence
Conjecture~\ref{conj:main}.
\end{proposition}

\begin{proof}
The root is legal, its envy graph being identically zero, and has
$\abs{R} = m$. If $m \le 3$ the statement is vacuous and
Conjecture~\ref{conj:cri} is a finite check; otherwise any legal assignment from
a state with $\abs{R} \ge 4$ lands in a state with $\abs{R} \ge 3$, which is
live by hypothesis, and a live state reaches a terminal by definition.
\end{proof}

This would have been a different shape of target from
Conjecture~\ref{conj:balance-rule}: that one asks for a global invariant on an
unbounded search, whereas Conjecture~\ref{conj:cri-depth} carries a constant, and
the constant is small. It is however false.

\begin{proposition}[The depth conjecture fails]
\label{prop:cri-depth-false}
Conjecture~\ref{conj:cri-depth} is refuted. There are reachable legal CR states
with $\abs{R} \ge 3$ from which no sequence of assignments and relabellings
reaches a terminal state through legal states.
\end{proposition}

\begin{proof}
Complete enumeration of the state space over instances drawn from four families
at $n = 3$ and $m \in \set{6,7,8}$ finds $657$ such states
(\texttt{updates/update\_48/depth\_stress.py}):
\begin{center}\small
\begin{tabular}{@{}lrrr@{}}
\toprule
family & $m = 6$ & $m = 7$ & $m = 8$ \\
\midrule
composed  & $7$  & $243$ & $137$ \\
capped    & $60$ & $210$ & $0$ \\
threshold & $0$  & $0$   & $0$ \\
nested    & $0$  & $0$   & $0$ \\
\bottomrule
\end{tabular}
\end{center}
Nothing is sampled within an instance: for each one every state is enumerated and
liveness is computed by a backward pass, so a reported dead state is dead.
\end{proof}

\begin{remark}[The caveat was the right one, and it fired]
\label{rem:cri-caveat}
Conjecture~\ref{conj:cri-depth} rested on $34{,}543$ states at $m \le 7$ with only
$6$ of them at $\abs{R} = 7$, and was recorded as a hypothesis rather than a
result on exactly that ground --- the lesson of \Cref{rem:n3-rules-fail}, where a
rule surviving $227$ instances failed on $368$. Two details of the refutation are
worth keeping. Every failure lies in the \emph{composed} family
$\cost_i(S) = f_i(\abs{S \cap D_i})$, where the residual of the solved cases
lives, or in \emph{capped}; threshold and nested produce none. And the capped
family had never been generated correctly before the defect of
\Cref{rem:cri-genbug} was repaired, so half of the refuting evidence was
unavailable while the conjecture was being formed.
\end{remark}

What Example~\ref{ex:cri-deadend} forces on any replacement is that its proof
cannot argue that the state stays good, since the state it must avoid is
optimal. It has to argue about what remains \emph{available}: by
Lemma~\ref{lem:cri-witness} every arc is computed on the contracted instance
$(D, \cost^{R})$, which by Lemma~\ref{lem:cri-contract} is again dichotomous, so
any such statement is internal to the class. \Cref{sec:cri-witness} takes that
route.

\begin{remark}[A defect in the test generators, and what it touched]
\label{rem:cri-genbug}
The adversarial families used here and in \Cref{sec:experiments} had a defect
found while checking Proposition~\ref{prop:cri-first}: in \texttt{f\_capped} and
\texttt{f\_threshold} the random parameter was drawn \emph{inside} the
comprehension over subsets, so it was redrawn for every $S$ and the resulting
functions were neither cappings nor thresholds --- they had marginals of $2$ and
more, and so lay \emph{outside} the dichotomous class. \texttt{f\_mixed} draws
from both and was affected too. The uniform, disjoint, nested and one-heavy
families were clean throughout.

The consequence is one of coverage rather than of correctness: those sweeps
tested a wider class than intended, so their negative findings stand while their
claimed coverage of capped and threshold instances does not. The
cycle-closing check of \Cref{rem:cyclebound-lossy} is unaffected, since
Theorem~\ref{thm:cyclebound} assumes only envy-freeability and not dichotomous
costs. With the generators repaired, the counterexample hunt still finds no
counterexample, and of the $29$ instances that come within one unit of being one,
$17$ are \emph{threshold} against $4$ nested and $3$ capped --- so
$\cost(S) = \max(0, \abs{S} - t)$ is the tightest family known, and it was not
previously being generated at all.
\end{remark}

% --------------------------------------------------------------------------
\subsection{Carrying a witness, and a bridge to Approach 10}
\label{sec:cri-witness}

Every invariant tried so far has been a \emph{predicate on the state}, and every
one has died: the three families eliminated in \Cref{sec:approach3}, and here
pointwise extendability, the last rung, balance, the commitment predicate and
Conjecture~\ref{conj:cri-depth}. Example~\ref{ex:cri-deadend} explains why, and
rules out the whole family at a stroke --- the state to be avoided there is
\emph{optimal}, with $\pathw{} = (0,0,0)$, so no measure of the current state's
quality can separate it from a live one. The standard response to that situation
is not another predicate but a stronger induction hypothesis: carry a
\emph{witness} alongside the state.

\begin{definition}[Completion]
\label{def:cri-completion}
A \emph{completion} of a CR state $(\alloc, R)$ is an allocation $\alloc^{\star}$
of all of $\items$ with $\bundle{i} \subseteq \alloc^{\star}_i$ for every $i$. It
is \emph{good} if $\pathw{\alloc^{\star}} \le 1$. Write
\[
  \Phi(\alloc,R) \;:\iff\; (\alloc,R) \text{ is legal and some completion of it
  is good.}
\]
\end{definition}

\begin{proposition}[The three easy clauses]
\label{prop:cri-phi}
$\Phi(\text{root})$ holds if and only if Conjecture~\ref{conj:main} holds for the
instance; at a terminal state $\Phi$ holds if and only if that state is good; and
$\Phi$ implies legality by definition.
\end{proposition}

\begin{proof}
The root has $\bundle{i} = \emptyset$ for every $i$, so every allocation of
$\items$ is a completion of it, and it is legal because its envy graph vanishes.
A terminal state has $R = \emptyset$, so it is its own unique completion.
\end{proof}

So, as always, progress is the whole content. What is new is that the standing
counterexample does not dispose of $\Phi$ before the work starts.

\begin{observation}[$\Phi$ survives Example~\ref{ex:cri-deadend}]
\label{obs:cri-phi-survives}
In Example~\ref{ex:cri-deadend} each of the three ways of placing the last chore
creates a positive-weight cycle, so that state has \emph{no} good completion and
$\Phi$ excludes it. Every predicate eliminated in \Cref{sec:approach3} and in
\Cref{sec:cri-refuted} was, by contrast, satisfied by some dead state.
\end{observation}

Pure witness-following is nevertheless not enough, and the reason is recorded
rather than guessed.

\begin{proposition}[Relabelling is unavoidable]
\label{prop:cri-follow-fails}
It is not the case that from every state satisfying $\Phi$ some chore may be
assigned to its owner in a good completion while preserving legality.
\end{proposition}

\begin{proof}
Such a rule would, by induction on $\abs{R}$ from the root, produce a schedule of
\emph{assignments alone} reaching a good terminal whenever
Conjecture~\ref{conj:main} holds for the instance. But three instances of the
complete exhaustive $n = m = 3$ family are bad roots in the assignment-only
frame while satisfying Conjecture~\ref{conj:main} and
Conjecture~\ref{conj:cri} (\Cref{obs:cri-sweep}).
\end{proof}

Hence a witness-carrying proof must let the target \emph{move}, which is exactly
what relabelling does: it permutes the bundles, and with them the completion
being carried.

\paragraph{The bridge.} Lemma~\ref{lem:cri-witness} says a CR state is a witness
on the contracted instance, and Lemma~\ref{lem:cri-contract} says that instance
is again dichotomous. Both hypotheses of \Cref{sec:approach10}'s main theorem are
therefore available, and in the CR frame relabelling \emph{is} the choice of
assignment, so a minimum-cost labelling may always be selected. Writing
\[
  \Sigma^{R}(\alloc) \;:=\; \sum_{i \in \agents}
  \Big[\max_t \cost^{R}_i(\bundle{t}) - \min_t \cost^{R}_i(\bundle{t})\Big]
\]
for the total spread of the family measured in the contracted costs:

\begin{theorem}[Bridge]
\label{thm:cri-bridge}
Let $n = 3$ and let $(\alloc, R)$ be a CR state with $\Sigma^{R}(\alloc) \le 3$.
Then every labelling of the family $\alloc$ minimising
$\sum_i \cost^{R}_i(\bundle{\sigma(i)})$ is legal.
\end{theorem}

\begin{proof}
By Lemma~\ref{lem:cri-contract} the contracted instance $(D, \cost^{R})$ is
negative dichotomous, and by Lemma~\ref{lem:cri-witness} the envy graph of the CR
state is the envy graph of the allocation $\alloc$ in that instance. Apply
Theorem~\ref{thm:f5-sigma3} to $(D,\cost^{R})$: a family of total spread at most
$3$ assigned by a minimum-cost matching is good, which is precisely legality of
the CR state.
\end{proof}

This is the first point at which the two lines meet: \Cref{sec:approach10}'s
theorem serves as the \emph{soundness} clause of this section's induction. The
base case is then immediate, and agrees with what was already proved.

\begin{corollary}[Base case]
\label{cor:cri-base}
From the root, assigning any chore $a$ to any agent yields
$\Sigma^{R'} = \sum_i \beta_i \le 3$, where
$\beta_i = \cost_i(\items) - \cost_i(\items \setminus \set{a})$; and the
minimum-cost labelling gives $a$ to an agent of minimal $\beta$, which is
Proposition~\ref{prop:cri-first}.
\end{corollary}

\begin{proof}
At the root every bundle is empty, so after the assignment
$\cost^{R'}_i(\set{a}) = \cost_i(\set{a} \cup R') - \cost_i(R') =
\cost_i(\items) - \cost_i(\items \setminus \set{a}) = \beta_i$ while the two
empty bundles have contracted cost $0$. Agent $i$'s three values are therefore
$(\beta_i, 0, 0)$ up to order, with spread $\beta_i \in \set{0,1}$, so
$\Sigma^{R'} = \sum_i \beta_i \le n = 3$. The cost of a labelling is $\beta_j$
for whichever agent $j$ receives $\set{a}$, minimised exactly at
$\beta_j = \min_k \beta_k$.
\end{proof}

\paragraph{The candidate invariant, and what is left.} Take
$\Phi_{\Sigma}(\alloc,R) :\iff \Sigma^{R}(\alloc) \le 3$. Its root clause is free
($\cost^{\items}_i \equiv 0$, so $\Sigma = 0$), its soundness clause is
Theorem~\ref{thm:cri-bridge}, and at a terminal $\cost^{\emptyset} = \cost$, so
$\Phi_\Sigma$ delivers a \emph{good allocation} --- that is,
Conjecture~\ref{conj:main} at $n = 3$. One clause remains.

\begin{remark}[The labelling must be got right, not waved through]
\label{rem:cri-labelling}
Every state \emph{visited} must be legal, so one may not assign and then relabel
to minimum cost: the un-relabelled state is visited on the way. A relabelling is
a single move between two states, both of which must be legal, so it cannot be
used to escape an illegal one. The correct form of the progress step therefore
quantifies over labellings as well as over chores.
\end{remark}

\begin{conjecture}[Progress --- \textbf{open}]
\label{conj:cri-progress}
Let $n = 3$ and let $\alloc$ be a family with $\Sigma^{R}(\alloc) \le 3$ and
$R \ne \emptyset$. Then there are a legal labelling of $\alloc$, a chore
$a \in R$ and a bundle of $\alloc$ such that adjoining $a$ to that bundle yields
a family $\alloc'$ with $\Sigma^{R'}(\alloc') \le 3$ whose induced labelling is
legal.
\end{conjecture}

\begin{proposition}
\label{prop:cri-progress-implies}
Conjecture~\ref{conj:cri-progress} implies Conjecture~\ref{conj:cri} at $n = 3$,
and hence Conjecture~\ref{conj:main} at $n = 3$.
\end{proposition}

\begin{proof}
The root satisfies $\Phi_\Sigma$ and is legal. Each application of
Conjecture~\ref{conj:cri-progress} decreases $\abs{R}$ by one through legal
states, so after $m$ steps a terminal state is reached with
$\Sigma^{\emptyset} \le 3$; Theorem~\ref{thm:cri-bridge} makes it good.
\end{proof}

\begin{remark}[What this route costs, and why the weaker target comes first]
\label{rem:cri-progress-cost}
At $\abs{R} = 1$ Conjecture~\ref{conj:cri-progress} produces a terminal family of
total spread at most $3$, so it implies Conjecture~\ref{conj:f5-sigma3}. The two
lines are therefore \emph{not} independent: this route to $n = 3$ contains
\Cref{sec:approach10}'s open statement. Two consequences, and they are the
practical content of this subsection. Conjecture~\ref{conj:f5-sigma3} is the
weaker target and should be settled first, since a refutation of it would refute
Conjecture~\ref{conj:cri-progress} without any further work. And if it is
refuted, $\Phi_\Sigma$ fails as an invariant and Conjecture~\ref{conj:cri} would
need a different one --- $\Phi$ of \Cref{def:cri-completion} being the candidate
that Observation~\ref{obs:cri-phi-survives} leaves standing.
\end{remark}

\paragraph{Verdict.} Conjecture~\ref{conj:cri} is open and unrefuted, and it is a
weaker target than either Conjecture~\ref{conj:h1prime} or
Conjecture~\ref{conj:balance-rule}. Lemmas~\ref{lem:cri-contract}
to~\ref{lem:cri-free}, Proposition~\ref{prop:cri-first} and
Theorem~\ref{thm:cri-bridge} are the first statement of this problem as an
induction on the \emph{instance} rather than a search over states, and they stand
whatever becomes of Conjecture~\ref{conj:cri-progress}. What the frame does
\emph{not} do is remove dead ends: it removes the ones of
Proposition~\ref{prop:deadends} and grows its own, and
Example~\ref{ex:cri-deadend} shows the new ones are harder to characterise, being
optimal states. Conjecture~\ref{conj:cri-depth} is refuted, and the surviving
route --- Conjecture~\ref{conj:cri-progress} --- is by
Remark~\ref{rem:cri-progress-cost} at least as strong as the open statement of
\Cref{sec:approach10}. As with every other approach in this document, no case
with $n \ge 3$ is proved here.
